%% ============================================================================
\section{The Run Mechanism}
\label{app:model}

The model below makes the run mechanism and the non-fundamental residual's friction term
explicit, and it yields an ex-ante cross-coin
prediction. The model does not produce the reported numbers. Section~\ref{sec:ident}'s
estimator derives those by accounting alone.

A unit mass of holders choose REDEEM or HOLD, with the HOLD payoff normalized to $0$ in
both states, so payoffs below are \emph{net gains of REDEEM over HOLD}. If the peg holds,
redeeming forfeits the holding benefit and pays the transaction cost, a net gain
\begin{equation}
g_n = -(\eta + c) < 0,
\end{equation}
where $\eta\ge 0$ is the forfeited yield and liquidity premium and $c\ge 0$ the redemption
cost. If the peg breaks, a holder who redeems early avoids the reserve haircut, a net gain
$g_b>0$. The realized liquid buffer is $\theta$, with $\mathbb{E}[\theta]\approx 1-\phi$,
and a run succeeds iff the redeeming mass $n$ reaches $\theta$.

By the global-games selection result
\cite{carlssonvandamme1993,morrisshin1998,goldsteinpauzner2005}, the model admits a unique
threshold equilibrium: the marginal holder at $\theta^*$ holds a uniform belief over the
redemption mass, $n\sim U[0,1]$, so $\Pr(\text{break})=\Pr(n\ge\theta^*)=1-\theta^*$.
Setting the REDEEM and HOLD payoffs equal gives $(1-\theta^*)\,g_b + \theta^*\,g_n = 0$.
With $g_n=-(\eta+c)$ this solves to
\begin{equation}
\theta^* = \dfrac{g_b}{g_b + \eta + c}.
\label{eq:thetastar}
\end{equation}
A self-fulfilling run occurs when $\theta<\theta^*$, so a larger $\theta^*$ means a
\emph{more} fragile coin.

The gain $g_b$ increases in a coin's reserve \emph{exposure} to the shock,
because more impaired backing means a larger haircut to avoid. By \eqref{eq:thetastar},
higher exposure implies a higher $\theta^*$ and greater fragility. This mechanism gives
the ex-ante fragility ordering applied in Section~\ref{sec:placebo}.

